% !PRINTED 26
\weberpage{26}

%<<CONT
vorkommt, und dass die Indices in diesem Theile der Reihe alle
gleich $\nu$ sind.

Die Behauptung ist richtig für $r = 2$, und wir beweisen sie
allgemein durch vollständige Induction, unter der Voraussetzung,
dass sie für $r$ schon als richtig erwiesen sei.

Ist $Q_r$ nicht gleich $R$, so wählen wir $Q_1^{(r)}$ so, dass $Q_{r+1}$ der
Durchschnitt von $Q_1, Q_1' \cdots Q_1^{(r-1)}, Q_1^{(r)}$ ist, und bezeichnen überdies
 den Durchschnitt von $Q_1, Q_1' \cdots Q_1^{(r-2)}, Q_1^{(r)}$ mit $Q_r'$. Dann
ist $Q_r'$ von $Q_r$ verschieden und $Q_{r+1}$ ist der Durchschnitt von
$Q_r$ und $Q_r'$. Nun haben wir nach der Voraussetzung (4) zwei
entsprechende Stücke der Compositionsreihe mit der constanten
Indexreihe $\nu$

\weq{5}{
\begin{array}{l}
Q,\ Q_1 \cdots Q_{r-1}\ Q_r,\\
Q,\ Q_1 \cdots Q_{r-1}\ Q_r';
\end{array}
}

denn die zweite dieser Reihen ist genau wie die erste nach der
Vorschrift (3) gebildet, nur dass $Q_1^{(r)}$ an Stelle von $Q_1^{(r-1)}$ getreten
ist. Nach dem Satze \S.\,6, 2. ist daher $Q_{r+1}$ grösster Normaltheiler
 sowohl von $Q_r$ als von $Q_r'$, und zwar vom Index $\nu$, und
es folgt also, was bewiesen werden sollte, dass die Compositionsreihe
 so fortgesetzt werden kann:

\weq{6}{Q,\ Q_1 \cdots Q_{r-1},\ Q_r,\ Q_{r+1},}

und dass der neu hinzugekommene Index gleichfalls $\nu$ ist.

Dies können wir nun fortsetzen, bis wir zur Gruppe $R$ gelangt
 sind, und erhalten ein Stück der Compositionsreihe

\weq{7}{
\begin{array}{ccc}
Q,\ Q_1,\ Q_2 & \cdots & R,\\
\nu\quad\ \nu & & \nu
\end{array}
}

deren letztes Glied $R$ wieder Normaltheiler von $P$ ist und in
der alle Indices übereinstimmen.

Setzen wir von $R$ an die Compositionsreihe fort, so kann
sich der Index ändern; er kann aber von da an wieder gleich
gehalten werden, bis man abermals zu einem Normaltheiler von
$P$ gelangt.

Damit ist der folgende Satz bewiesen:

III.\ \gesperrt{Man kann die Compositionsreihe von $P$ mit ihrer
Indexreihe}

\weq{8}{
\begin{array}{c}
P,\ P_1,\ P_2,\ P_3 \cdots\\
j_1,\ j_2,\ j_3 \cdots
\end{array}
}

\gesperrt{so anordnen, dass, so oft eine Aenderung in der
Indexreihe eintritt, also $j_{r+1}$ von $j_r$ verschieden
ist, $P_r$ ein Normaltheiler von $P$ ist.}
